\documentclass[aspectratio=169,10pt]{beamer}
\usetheme[numbering=fraction,progressbar=frametitle]{metropolis}
\usepackage{fontspec}
\setsansfont{Latin Modern Sans}
\usepackage{amsmath,amssymb,mathtools,graphicx,booktabs}
\usepackage{polyglossia}
\setdefaultlanguage{french}
\setbeamertemplate{navigation symbols}{}
\setbeamerfont{frametitle}{size=\large}
\setbeamerfont{normal text}{size=\normalsize}
\setlength{\parskip}{.45em}
\hypersetup{pdfauthor={},pdfsubject={Lecture slides},pdfkeywords={},pdftitle={Statistiques}}
\providecommand{\N}{\mathbb{N}}
\providecommand{\Z}{\mathbb{Z}}
\providecommand{\D}{\mathbb{D}}
\providecommand{\Q}{\mathbb{Q}}
\providecommand{\R}{\mathbb{R}}
\providecommand{\C}{\mathbb{C}}
\providecommand{\vect}[1]{\overrightarrow{#1}}
\providecommand{\vi}{\vec{\imath}}
\providecommand{\vj}{\vec{\jmath}}
\providecommand{\vk}{\vec{k}}
\providecommand{\norm}[1]{\left\lVert #1\right\rVert}
\providecommand{\vc}[2]{\begin{pmatrix}#1\\#2\end{pmatrix}}
\providecommand{\deter}[2]{\det\left(#1,#2\right)}
\providecommand{\abs}[1]{\left\lvert #1\right\rvert}
\providecommand{\dd}{\,\mathrm{d}}
\providecommand{\e}{\mathrm{e}}
\providecommand{\Card}{\operatorname{Card}}
\providecommand{\E}{\operatorname{E}}
\providecommand{\V}{\operatorname{V}}
\providecommand{\degre}{^{\circ}}
\providecommand{\un}[1]{\,\mathrm{#1}}
\providecommand{\sci}[2]{#1\times10^{#2}}
\providecommand{\pgcd}{\operatorname{PGCD}}
\providecommand{\ppcm}{\operatorname{PPCM}}
\title{Statistiques}
\subtitle{Cours et exemples}
\date{}
\author{}
\begin{document}
\begin{frame}[plain]
\vfill
{\usebeamerfont{title}\inserttitle\par}
\medskip
{\usebeamerfont{subtitle}\insertsubtitle\par}
\medskip
{\usebeamercolor[fg]{progress bar}\rule{\textwidth}{.5pt}}
\vfill
\end{frame}
\begin{frame}{}
Comparer les séries \(A=(8,8,8,8,8)\) et \(B=(4,6,8,10,12)\). Elles ont la même
moyenne. Sont-elles pourtant semblables? Proposer une mesure de la dispersion,
puis confronter les propositions à la variance et à l'écart-type.
\end{frame}
\begin{frame}{Objectifs}
\begin{itemize}

  \item Organiser des données statistiques : variable discrète et continue.
  
\item Calculer et interpréter les indicateurs de position :
        moyenne, médiane, mode.
  
\item Calculer et interpréter les indicateurs de dispersion :
        étendue, variance, écart-type.
  
\item Construire et lire les représentations graphiques :
        diagramme en bâtons, diagramme circulaire, histogramme.

\end{itemize}
\end{frame}
\begin{frame}{Statistiques descriptives}
Pour une série non vide de valeurs \(x_1,\ldots,x_n\) :

\begin{itemize}

  \item \textbf{Moyenne} : \(\bar{x} = \dfrac{1}{n}\displaystyle\sum_{i=1}^n x_i\).
  
\item \textbf{Médiane} \(Q_2\) : valeur centrale une fois les données triées.
  
\item \textbf{Quartiles} \(Q_1\), \(Q_3\)  : valeurs aux rangs \(\lceil n/4\rceil\) et \(\lceil3n/4\rceil\) dans la série triée.
  
\item \textbf{Écart-type} :
        \(\sigma = \sqrt{\dfrac{1}{n}\displaystyle\sum_{i=1}^n (x_i - \bar{x})^2}\).

\end{itemize}
\end{frame}
\begin{frame}{Caractéristiques complémentaires}
Le \textbf{mode} d'une série statistique est la valeur (ou la classe) qui
apparaît le plus souvent. Dans une distribution en classes, on parle de
\textbf{classe modale}.
\end{frame}
\begin{frame}{Caractéristiques complémentaires}
La \textbf{fréquence} d'une valeur \(x_i\) est \(f_i = n_i / N\).
La \textbf{fréquence cumulée croissante} jusqu'à \(x_k\) est
\(F_k = \displaystyle\sum_{i \leq k} f_i\).
La médiane se détermine par les rangs centraux dans la série triée.
La fréquence cumulée n'atteint pas forcément exactement \(50\%\).
\end{frame}
\begin{frame}{Séries statistiques groupées}
Pour \(n\) valeurs \(x_1,\ldots,x_p\) aux effectifs \(n_1,\ldots,n_p\)
(avec \(N=\sum n_i\)), la \textbf{fréquence} de \(x_i\) est \(f_i=n_i/N\).
\end{frame}
\begin{frame}{Séries statistiques groupées}
\begin{itemize}

  \item \textbf{Moyenne :} \(\bar{x} = \dfrac{1}{N}\sum_{i=1}^p n_i x_i\).
  
\item \textbf{Médiane :} valeur \(Q_2\) partageant la série en deux moitiés égales.
  
\item \textbf{Mode :} valeur (ou classe) de plus grande fréquence.
  
\item \textbf{Quartiles :} \(Q_1\) (25\%), \(Q_2\) (50\%), \(Q_3\) (75\%).

\end{itemize}
\end{frame}
\begin{frame}{Séries statistiques groupées}
\begin{itemize}

  \item \textbf{Étendue :} \(e = x_{\max} - x_{\min}\).
  
\item \textbf{Variance :}
        \(V = \dfrac{1}{N}\sum_{i=1}^p n_i(x_i-\bar{x})^2
           = \overline{x^2} - \bar{x}^2\)
        (variance \(=\) moyenne des carrés moins carré de la moyenne).
  
\item \textbf{Écart-type :} \(\sigma = \sqrt{V}\).

\end{itemize}
\end{frame}
\begin{frame}{Séries statistiques groupées}
\begin{enumerate}
\setcounter{enumi}{0}
  \item Trier les valeurs dans l'ordre croissant.
  
\item Calculer les effectifs cumulés.
  
\item La médiane est :

\begin{itemize}

      \item si \(N\) est impair : la valeur de rang \((N+1)/2\).
      
\item si \(N\) est pair : la demi-somme des valeurs de rangs \(N/2\) et \(N/2+1\).

\end{itemize}

\end{enumerate}
\end{frame}
\begin{frame}{Séries statistiques groupées}
Pour des données connues seulement par classes, remplacer les valeurs par
les centres donne une \textbf{estimation} de la moyenne et de la variance.
La médiane exacte n'est pas déterminée par les seuls effectifs des classes.
Les estimations par interpolation supposent une répartition uniforme dans
chaque classe. Dans un histogramme, ce sont les \textbf{aires} des rectangles
qui sont proportionnelles aux effectifs; pour des amplitudes inégales,
les hauteurs sont proportionnelles à effectif/amplitude.
\end{frame}
\end{document}
